\section{Introdução}
Queremos definir, dentro da teoria dos conjuntos, o que significa um conjunto ter $0$, $1$, $2$, $3$, \dots\ elementos.

Para isso, precisamos escolher conjuntos específicos que possam representar os números naturais.

A ideia central é a seguinte: dois conjuntos devem representar o mesmo número natural quando têm a mesma quantidade de elementos. Isso nos leva ao conceito de equipotência.

Mesmo fora da matemática, é possível comparar quantidades sem recorrer previamente aos números.
Para registrar a quantidade de ovelhas em um pasto, pode-se levar um saco de pedras e, para cada ovelha contada, colocar uma pedra no saco.
Mais tarde, ao comparar as ovelhas com as pedras, verifica-se se a correspondência é exata.

\begin{definition}
    Sejam $A$, $B$ conjuntos.
    Dizemos que $A$ e $B$ possuem o mesmo número de elementos (ou que são equipotentes) se existe uma função bijetora de $A$ em $B$.
    Escrevemos $A \approx B$ para denotar que $A$ e $B$ são equipotentes.
    
    Formalmente, $\forall A \forall B\, (A \approx B \leftrightarrow \exists f (f:A\to B \land f \text{ é bijetora}))$.
\end{definition}

Seguem as propriedades mais elementares da equipotência.
\begin{proposition}
    A equipotência é uma relação de equivalência no universo dos conjuntos.
    Mais formalmente:
    \begin{enumerate}[label=(\alph*)]
        \item $\forall A \, A\approx A$.
        \item $\forall A \forall B\, (A \approx B \rightarrow B \approx A)$.
        \item $\forall A \forall B \forall C\, (A \approx B \land B \approx C \rightarrow A \approx C)$.
    \end{enumerate}
\end{proposition}
\begin{proof}
    Para (a), seja $A$ um conjunto.
    A função identidade em $A$, dada por $\id_A = \{(x, x) : x \in A\}$, é uma função bijetora de $A$ em $A$, e, portanto, $A \approx A$.

    Para (b), sejam $A$ e $B$ conjuntos tais que $A \approx B$.
    Então existe uma função bijetora de $A$ em $B$, digamos $f:A\to B$.
    Segue que $f^{-1}:B\to A$ é bijetora, e, portanto, $B \approx A$.

    Para (c), sejam $A$, $B$ e $C$ conjuntos tais que $A \approx B$ e $B \approx C$.
    Então existe uma função bijetora de $A$ em $B$, digamos $f:A\to B$, e existe uma função bijetora de $B$ em $C$, digamos $g:B\to C$.
    Assim, a composição de funções é uma função bijetora de $A$ em $C$, ou seja, existe uma função bijetora de $A$ em $C$, e, portanto, segue que $A \approx C$.
\end{proof}

Em geometria plana, é relativamente comum definir geometricamente um vetor como a classe de equivalência de segmentos orientados, onde dois segmentos orientados são equivalentes se são paralelos, possuem o mesmo comprimento e a mesma direção.

Uma primeira ideia para definir-se, digamos, o número $1$, seria adotar uma estratégia análoga e definir o número $1$ como a classe de equivalência de conjuntos que tem exatamente um elemento.
Ou seja, o ``conjunto de todos os singletos''.
De forma análoga, o número $2$ seria o ``conjunto de todos os pares não ordenados que não são singletos'', e ``assim por diante''.
Não está imediatamente claro o que o ``e assim por diante'' significaria, mas não divagaremos sobre isso, uma vez que o lema abaixo nos mostra que essa estratégia é invariavelmente condenada ao fracasso.
\begin{lemma}
    Não existe um conjunto $A$ tal que, para todo conjunto $x$, $\{x\} \in A$.
\end{lemma}

\begin{proof}
    Suponha que exista tal conjunto $A$ e considere $V=\bigcup A$.

    Mostraremos que todo conjunto pertence a $V$.
    Seja $x$ um conjunto arbitrário.
    Pela hipótese, $\{x\}\in A$.
    Como $x\in\{x\}$, segue da definição de união que $x\in\bigcup A=V$.

    Portanto, $\forall x\, (x\in V)$.
    Mas já foi mostrado anteriormente que não existe um conjunto com essa propriedade.
    Chegamos, assim, a uma contradição.
\end{proof}

Seguimos para a estratégia seguinte.
Ao invés de definirmos o número $1$ como a classe de equivalência de todos os singletos, definiremos o número $1$ como um \emph{representante} dessa classe.

Mas qual escolher?
Podemos tomar $\{\{\{\}\}\}$, ou $\{\{\}\}$, ou $\{\{\{\{\}\}\}\}$, e, em um primeiro momento, não parece haver nenhuma razão para escolher um representante em detrimento de outro.

Vamos tentar racionalizar um processo que torne tal escolha natural.
Comecemos com o número $0$.
Queremos que $0$ corresponda a um conjunto sem nenhum elemento.
Ora, tal escolha é trivial, pois só há um tal conjunto, o vazio.

Define-se, assim, $0=\emptyset$.

O próximo passo é definir o número $1$.
Queremos um conjunto que tenha exatamente um elemento, ou seja, um singleto.
Para tanto, precisamos escolher um tal elemento.
Ora, já temos um conjunto em mãos, o $0$. Assim, define-se $1=\{0\}$.

Agora já escolhemos conjuntos correspondentes ao $0$ e ao $1$.
O próximo passo é definir o número $2$.
Queremos um conjunto que tenha exatamente dois elementos, ou seja, um par não ordenado. Para tanto, precisamos escolher dois tais elementos.
Ora, já temos dois conjuntos em mãos, o $0$ e o $1$. Assim, define-se $2=\{0,1\}$.

De forma geral, queremos definir:

\begin{align*}
    0 &= \emptyset,\\
    1 &= \{0\} &= 0 \cup \{0\},\\
    2 &= \{0,1\} &= 1 \cup \{1\},\\
    3 &= \{0,1,2\} &= 2 \cup \{2\},\\
     &\vdots
\end{align*}

Porém, tal definição é informal, devido às reticências.

Porém, parece que encontramos um mecanismo de ``escolher o próximo número'', que é formalizado pela definição a seguir.

\begin{definition}
    Dado um conjunto $x$, o \emph{sucessor} de $x$ é o conjunto $S(x) = x \cup \{x\}$.
\end{definition}

Com isso em mente, como definir o conjunto dos números naturais?
Ora, tal conjunto deve ter o $0$, e, sempre que tem um conjunto $x$, deve ter seu sucessor $S(x)$.
Um conjunto com essa propriedade é chamado de conjunto indutivo.

\begin{definition}
    Dizemos que um conjunto $A$ é \textbf{indutivo} se $0 \in A$ e, para todo $x \in A$, $S(x) \in A$.
\end{definition}

Com os axiomas que temos até então não é possível garantir a existência de um conjunto indutivo.
O Axioma do Infinito diz precisamente que um tal conjunto existe.

\begin{axiom}
    Existe um conjunto indutivo.
    Formalmente,
    \[
        \exists I\,\bigl(0\in I \land \forall x\,(x\in I \rightarrow S(x)\in I)\bigr).
    \]

    Escrevendo apenas na linguagem original, tal axioma se traduz como:

    \[
    \exists I\,\Bigl(
        \exists z\,(\forall w\,(w\notin z) \land z\in I)
        \land
        \forall x\,(x\in I \rightarrow \exists y\,(y\in I \land \forall z\,(z\in y \leftrightarrow (z\in x \lor z=x))))
    \Bigr).
    \]
\end{axiom}

Um conjunto indutivo $I$ tem $0$ como elemento por definição.
Portanto, também tem $S(0)=\{0\}=1$ como elemento.
Assim, $I$ tem $S(1)=\{0,1\}=2$ como elemento.
Intuitivamente, ele deverá conter todos os números naturais.
Porém, nada garante que ele não contenha outros elementos além destes.

Buscaremos então definir o \emph{menor conjunto indutivo}, ou seja, um conjunto indutivo que está contido em todos os outros, mínimo com relação à inclusão $\subseteq$.

\begin{definition}
    O conjunto dos números naturais, denotado por $\omega$, é o único conjunto $J$ tal que
    \[
        \forall x\, \bigl(x \in J \leftrightarrow \forall I\, (I \text{ é indutivo} \rightarrow x \in I)\bigr).
    \]
\end{definition}

Vamos mostrar que tal conjunto existe e é único.

\begin{lemma}
    Existe um único conjunto $J$ como acima.
\end{lemma}
\begin{proof}
    Seja $I'$ um conjunto indutivo, o que existe pelo Axioma do Infinito.

    Seja $J=\{x \in I' : \forall I\, (I \text{ é indutivo} \rightarrow x \in I)\}$.

    Mostraremos que $J$ tem a propriedade desejada.

    Seja $K$ um conjunto indutivo.
    Segue que $J\subseteq K$, pois, para todo $x\in J$, $x$ é tal que $\forall I\, (I \text{ é indutivo} \rightarrow x \in I)$, e, portanto, $x\in K$.

    Reciprocamente, seja $x$ tal que $\forall I\, (I \text{ é indutivo} \rightarrow x \in I)$.
    Como $I'$ é indutivo, segue que $x\in I'$.
    Assim, $x\in J$.

    Para a unicidade, sejam $J_1$ e $J_2$ conjuntos com a propriedade desejada.
    Segue que para todo $x$:

    \[
        x \in J_1 \leftrightarrow \forall I\, (I \text{ é indutivo} \rightarrow x \in I) \leftrightarrow x \in J_2.
    \]

    Assim, pelo Axioma da Extensão, $J_1=J_2$.
\end{proof}

Vejamos algumas propriedades básicas de $\omega$.

\begin{proposition}
    $\omega$ é um conjunto indutivo.
\end{proposition}
\begin{proof}
    $0\in\omega$, pois, para todo conjunto indutivo $I$, $0\in I$.

    Seja $x\in\omega$.
    Assim, $\forall I\, (I \text{ é indutivo} \rightarrow x \in I)$.
    Seja $I$ um conjunto indutivo.
    Segue que $S(x)\in I$, pois $I$ é indutivo.
    Assim, $\forall I\, (I \text{ é indutivo} \rightarrow S(x) \in I)$, e, portanto, $S(x)\in\omega$.
\end{proof}

\begin{proposition}
    O único subconjunto indutivo de $\omega$ é $\omega$.
\end{proposition}

\begin{proof}
    Temos que $\omega$ é um conjunto indutivo.
    Se $J\subseteq \omega$ é indutivo, segue da definição de $\omega$ que $\omega \subseteq J$.
    Assim, $J=\omega$.
\end{proof}

Agora é possível demonstrar que $\omega$ satisfaz o esquema da indução.

\begin{proposition}
    Considere uma fórmula $\phi(x, t_1, \dots, t_n)$.

    O seguinte é um teorema:
    \[
    \forall t_1, \dots, t_n \Bigl(
        \bigl(
            \phi(0, t_1, \dots, t_n)
            \land
            \forall x \in \omega\, (\phi(x, t_1, \dots, t_n) \rightarrow \phi(S(x), t_1, \dots, t_n))
        \bigr)
        \rightarrow
        \forall x \in \omega\, \phi(x, t_1, \dots, t_n)
    \Bigr).
    \]
\end{proposition}

\begin{proof}
    Sejam dados $t_1, \dots, t_n$.
    Considere o conjunto $I = \{x \in \omega : \phi(x, t_1, \dots, t_n)\}$.
    Mostraremos que $I$ é indutivo.

    Por hipótese, $\phi(0, t_1, \dots, t_n)$, e, portanto, $0\in I$.
    Seja $x\in I$.
    Assim, $\phi(x, t_1, \dots, t_n)$ e $x \in \omega$.
    Como $\omega$ é indutivo, segue ainda que $S(x)\in\omega$.
    Por hipótese, $\phi(S(x), t_1, \dots, t_n)$, e, portanto, $S(x)\in I$.
    Assim, $I$ é indutivo.

    Como o único subconjunto indutivo de $\omega$ é $\omega$, segue que $I=\omega$.
    Assim, $\forall x \in \omega \, \phi(x, t_1, \dots, t_n)$.
\end{proof}

Podemos agora formalmente definir os $10$ primeiros números naturais:
\begin{definition}
    \begin{align*}
        0 &= \emptyset,\\
        1 &= S(0) = \{0\},\\
        2 &= S(1) = \{0,1\},\\
        3 &= S(2) = \{0,1,2\},\\
        4 &= S(3) = \{0,1,2,3\},\\
        5 &= S(4) = \{0,1,2,3,4\},\\
        6 &= S(5) = \{0,1,2,3,4,5\},\\
        7 &= S(6) = \{0,1,2,3,4,5,6\},\\
        8 &= S(7) = \{0,1,2,3,4,5,6,7\},\\
        9 &= S(8) = \{0,1,2,3,4,5,6,7,8\}.
    \end{align*}
\end{definition}